\begin{document}
	\maketitle
	\vspace{-6em} %设置摘要与标题的间距
	\zihao{-4} %设置正文字号
	%摘要部分
	\begin{abstract}
%<paper-block id="abstract.p001" type="paragraph">
针对测向误差有界、源数量及接收半径未知的干扰源定位与清除任务，本文将交会几何、接收安全选点和任务时间决策统一组织为四个递进模型。
%</paper-block>

%<paper-block id="abstract.p002" type="paragraph">
\textbf{针对问题一}，将每条带 $\pm1^\circ$ 误差的示向度表示为两个半平面，区分空集、无界和有界交会状态；对有界区域利用最远顶点对计算连续直径，给出直径圆覆盖的充要条件，并以两次测向与等边三角形反例说明覆盖不总成立，得到 $1\le2R^*/D\le2/\sqrt3$。
%</paper-block>

%<paper-block id="abstract.p003" type="paragraph">
\textbf{针对问题二}，以保证第二次接收或近距返回、最小化最坏后验直径为目标，推导四圆接收安全域，通过全域下界和整数区间证书的相等上界，得到标准未截断状态的两个最优站位 $(843.034754,\pm545.528422)$ 米，最坏直径为 $110.969319$ 米。通用规则达到策略级极小极大保证，但不宣称每个边界截断状态都单独最优；1500个配对场景中均取得第二次方向，平均后验直径上界为50.64米。
%</paper-block>

%<paper-block id="abstract.p004" type="paragraph">
\textbf{针对问题三}，利用全向源的正、负观测维护定位区与逐频道排空证据，将补扫点插入开放路径，并按移动、检测、切频道及光学清除的总成本选择动作。有限光学覆盖与公开的16源数量上界共同形成清除及停止判据。8场同版本演练共98个源全部清除；3次正式测试分别清除11、12、16个源并正常结束，按已清除源数加权的平均时间为233.01秒/源。
%</paper-block>

%<paper-block id="abstract.p005" type="paragraph">
\textbf{针对问题四}，不再把无信号解释为距离排除，以共同近邻凸包条件和连续单元验证建立21站发现方案，交错安排扫描与清除，并在已有移动线段上插入有预算的补测。在168个离线确认及边界场景中均全清且正常停止，平均每源时间较基础搜索策略下降2.85\%；3场当前版本演练全清，3条正常正式入口记录共清除40个源，加权时间约420.94秒/源。有限测试支持平均成本改善，不代表逐场占优或全局时间最优；正式记录的内部完成判定不冒充官方公开真值全清率。
%</paper-block>

%<paper-block id="abstract.p006" type="paragraph">
\noindent\textbf{关键词：} 有界测向误差；交会定位；极小极大；连续覆盖；搜索清除；路径规划
%</paper-block>
\end{abstract}
		
		\clearpage %换页
		
		%正文部分
		%Part one
